\section{Synthesiser estimator}
\label{app:synth}

The cross-arm analysis using only majority vote is given in Appendix~\ref{app:cross-arm}. That
analysis is cleaner --- no training is used and any difference between arms is due to the prompt
difference only (the synthesiser used a separate trained model for each arm, inflating its variance
with training variance in addition to the prompt). The majority-only analysis is relegated to the
Appendix because it does not alter any of our conclusions.

For the synthesiser accuracies on 184 test problems, we get 0.783 (op), 0.804 (pl), and 0.799
(neutral). None of the contrasts between any two arms excludes zero: the op--pl contrast is $-0.022$
($[-0.065, +0.022]$, $q = 0.85$), the op--neutral contrast is $-0.022$ ($[-0.065, +0.027]$,
$q = 0.85$), and the pl--neutral contrast is $+0.005$ ($[-0.027, +0.038]$, $q = 0.88$). The
point-estimates rank the operation arm lower than the controls, with the opposite ranking as with
majority voting, but the $\pm 0.05$ size of the confidence intervals at $n = 184$ exclude any
ranking.

The operation arm's synthesiser accuracy is not significantly different from the baselines it is
designed to outperform: there is no difference to pure majority voting ($\Delta = 0.000$) and the
biggest difference between the synthesiser and any one of the six single agents is 0.027 in absolute
value, with the confidence interval including zero. From \texttt{runs/cross\_arm.json} and
\texttt{runs/arms-cca\_council/evaluation.json}.

\section{Provenance and exclusion accounting}
\label{app:provenance}

We leave provenance to be inferred: providers self-report whether their answers are fabricated
(omitted answers are treated as synthetic to guarantee that a buggy provider logs simulated, not
real, records; if any provider is synthetic, the record is synthetic; and the analysis step simply
skips simulated records); track and report refusals, truncation, malformed outputs, and retried
queries by role (rather than assuming them); and omit rows without a record in any variant from
contrasts by role (rather than assuming them to be zero).

We report no refusals, no truncations, no malformed outputs, and no retried queries in our current
experiments: we encountered none in 1,128 pilot queries or 3,600 queries in the manipulation check
study.

\section{Scope of the novelty search}
\label{app:novelty-search}

We only used nine queries to arXiv, ACL Anthology, OpenReview (rejected and withdrawn papers), and
Semantic Scholar and only performed title/abstract screening, did not obtain full-texts nor snowball
through citations, nor queried PubMed by clinical MeSH terms, nor even had a double screener to
report inter-screener agreement on, reported only five out of eight queries to Semantic Scholar that
encountered rate limits as failures (rather than as zero findings), already encountered two
potential predecessors and one borderline finding (\S\ref{sec:related}) and thus claim no ``none in
$N$ papers screened''.

\section{What the check design concedes}
\label{app:concedes}

Three concessions were predeclared; a fourth was identified by post-collection audit and is reported
here.

\paragraph{The shared output contract compresses the contrast.}
To all versions we explicitly request atomic claims, counterclaims, and fault indicators in their
output. This means that even the length-matched neutral condition will output some counterclaims. In
that case we can only check for quantity, not for existence. While the dynamic range persists, it
will be squeezed compared to a setting where the neutral version did not output such information. We
anticipate operation $-$ neutral gaps to be smaller than with a version of this experiment that did
not employ a common record format but viewed the standard record format as necessary to route all
treatments through the same measurement apparatus.

\paragraph{Absent records are excluded, never scored as zero.}
We do not impute missing values, i.e.\ in case of refusals, truncations, or other malformations.
Otherwise we would artificially deflate the check scores in line with the failure rate, rather than
the performance of the variations.

\paragraph{The checks are benchmark-independent.}
A check only evaluates if the model lived up to its role, but does not evaluate the utility of the
role. For this reason shrinking the scope from four to two benchmarks is cost-neutral: it comes at
the expense of losing a role-wise distinction on a designated primary benchmark, not any roles.

\paragraph{The local-claim guard collapsed, and the check is therefore one-factorial in fact.}
\texttt{\_check\_claim\_decomposition} scores atomic claims per record weighted by the fraction
carrying a local check, so the weighting term is the guard: a role that emitted many unchecked
claims should not score well. In the collected records that term is the constant 2.0 in 1,191 of
1,200 records and has no variance at all in the operation arm of either family
(Table~\ref{tab:localclaim-guard}). The check is thus $2\times$ the claim count, and its guard never
discriminated. We report all local-claim results for preregistration transparency, but do not
interpret them as a validated pass or reversal.

\begin{table}[h]
\centering
\small
\begin{tabular}{llrrrl}
\toprule
Family & Arm & $n$ & weight $= 2.0$ & weight $< 2.0$ & distinct values \\
\midrule
\texttt{gpt-5.6-terra} & operation & 200 & 200 & 0 & \{2.0\} \\
\texttt{gpt-5.6-terra} & persona   & 200 & 200 & 0 & \{2.0\} \\
\texttt{gpt-5.6-terra} & neutral   & 200 & 200 & 0 & \{2.0\} \\
\texttt{grok-4.3}      & operation & 200 & 200 & 0 & \{2.0\} \\
\texttt{grok-4.3}      & persona   & 200 & 197 & 3 & \{1.667, 2.0\} \\
\texttt{grok-4.3}      & neutral   & 200 & 194 & 6 & \{1.667, 2.0\} \\
\bottomrule
\end{tabular}
\caption{The local-claim check's source-weighting guard, by family and arm. The guard takes one value
in 1,191 of 1,200 records and is constant in every operation arm, so the product it belongs to has
only one live factor.}
\label{tab:localclaim-guard}
\end{table}

\paragraph{Two further guards are defeatable rather than collapsed.}
The consistency safeguard evaluates flag rate against abstain rate. This means that the behavior of
always-flagging but never-abstaining gets a perfect score of 1.0 (compared to the operation prompt
score of 0.025). This safeguards against always-flagging-and-abstaining, but not
always-flagging-and-responding. The calibration safeguard was designed with the assumption that
uniformly reducing confidence would not reduce the expected calibration error. However, in this
dataset, we see that uniformly reducing confidence by a constant for all samples reduces ECE from
0.144 to 0.112 ($-0.05$ on the operation arm), and from 0.129 to 0.065 ($-0.10$ on the neutral arm),
without actually performing any calibration based on individual samples. While neither of these were
pre-registered, they do put limitations on how you can interpret their corresponding safeguard
metric in \S\ref{sec:did-fire}, but they do not supplant the safeguard metric.

\section{Benchmark materialisation and splits}
\label{app:splits}

We split the data 50/15/15/20 train/calibration/validation/test by calculating the salted hash of the
item contents. We further sample 200 items from the train split to be used for a manipulation-check
dev set.

We draw 1,000 HaluEval samples per arm for the benchmark arms, sampled sequentially from the file.
Sampling sequentially is not correlated with the data split (which is calculated off of item
content, not file index), and we have confirmed that the proportions of items in each split of the
subsample are within 1 percentage point of those in the full corpus.

\section{Failure modes in coalition attribution}
\label{app:attribution-failures}

If manipulation checks are an assessment of whether each role followed its instructions, this section
evaluates whether the attribution system followed its instructions: did the attribution system
calculate that a role's contribution to a coalition ablation was $\phi$, or does it vouch for
hallucination? We present two ways that this fails, discovered by running code that we had written
but not executed. Two are from our system presented as existence proofs with explanations of
mechanism, not prevalence in the wild.

\subsection{Finding (i): an all-zero ablation that passed its own efficiency check}

When we ran a Shapley ablation in retrain-per-coalition mode, it completed all 64 coalitions without
error and returned identically zero contributions for every role, while the efficiency axiom --- the
standard internal correctness check for a Shapley computation --- reported no violation.

\begin{table}[h]
\centering
\small
\begin{tabular}{lrr}
\toprule
 & retrain & masked (same config) \\
\midrule
Distinct coalition values across 64 & 1 & 30 \\
$v(\emptyset)$ & 0.0000 & 0.0000 \\
$v(A)$ & 0.0000 & 0.3314 \\
Every role's $\phi$ & 0.0000, CI $[0, 0]$ & $+0.036 \ldots +0.084$ \\
Efficiency check & 0.00e-00 (vacuous) & 0.00e-00 \\
\bottomrule
\end{tabular}
\caption{The degenerate ablation beside a masked run on identical configuration. The efficiency check
passes in both columns, and distinguishes them not at all.}
\label{tab:degenerate}
\end{table}

Why? We also found that we were reusing the same calibration data for all models. We had calibrated
the logit temperature and selective prediction threshold on the calibration set for the original
model trained on the main coalition dropout, and then reused this for all 64 models that are
independently retrained for each new data dropout. This meant the threshold was set way higher than
all selection scores, which are a completely new scale for each model, causing all to abstain and all
coalition values to be equal to the abstention credit (i.e., zero).

No overfitting was observed since efficiency (i.e., $\sum_i \phi_i = v(A) - v(\emptyset)$) will
always hold for $0 = 0 - 0$. Anything true for a non-degenerate answer is also true for a degenerate
answer; \textbf{this is the key learning, as internal checks can never catch a degenerate solution
because degenerate solutions are always internally consistent.}

The solution to this was to train each model in the ablation loop with its own logit temperature and
selective prediction threshold, both calibrated solely on the calibration set. Generalised learning:
when ablation involves retraining of a model, recalibrate all models. This is not a detail you can
cut corners on without completely breaking your analysis.

This guard is a different efficiency check --- a degeneracy check tallying distinct values assigned
to coalitions and loudly failing if the value function is degenerate. And indeed our scoped
experiment on real records now checks out, reporting 5 distinct coalition values:
\texttt{OK got 5 distinct coalition values}, a sanity check that the patched branch works correctly
on real data.

\subsection{Finding (ii): the masking approximation exceeds the effect it is used to detect}

With the calibration fixed, both estimator modes were run on identical configuration, differing only
in \texttt{shapley\_mode}. This isolates the masking approximation.

\begin{table}[h]
\centering
\small
\begin{tabular}{lrrr}
\toprule
Role & masked & retrain & gap \\
\midrule
\texttt{counter\_evidence\_checker} & $+0.0618$ & $+0.0356$ & $-0.0262$ \\
\texttt{consistency\_checker}       & $+0.0381$ & $+0.0510$ & $+0.0129$ \\
\texttt{local\_claim\_checker}      & $+0.0566$ & $+0.0711$ & $+0.0145$ \\
\texttt{explore\_then\_prune}       & $+0.0836$ & $+0.0607$ & $-0.0230$ \\
\texttt{rare\_hypothesis\_checker}  & $+0.0556$ & $+0.0560$ & $+0.0004$ \\
\texttt{confidence\_calibrator}     & $+0.0356$ & $+0.0570$ & $+0.0214$ \\
\midrule
\multicolumn{4}{l}{mean $|$gap$|$ 0.0164 \quad max $|$gap$|$ 0.0262 \quad SESOI 0.010} \\
\bottomrule
\end{tabular}
\caption{Per-role contributions under the two estimator modes. The mean disagreement between
estimators is larger than the smallest effect the pre-registration declares worth detecting.}
\label{tab:estimator-gap}
\end{table}

Two consequences follow, and both are properties of the estimator rather than of the data.

\textbf{The approximation error exceeds the target effect.} For per-role effect sizes smaller than
0.010 (the lowest effect size of interest in the pre-registration), the mean gap between the two
modes is 0.0164, so a masked experiment will not be able to report a per-role conclusion at the
effect size of interest in the pre-registration, regardless of the confidence interval.

\textbf{The ranks are not stable.} The \texttt{explore\_then\_prune} role ranks 1st and the
\texttt{counter\_evidence\_checker} role ranks 2nd according to the masked ranks, but after retrain,
the \texttt{local\_claim\_checker} ranks 1st and the \texttt{counter\_evidence\_checker} ranks last.
Statements like ``role X contributes the most'' are an artifact of the estimator and must be
calculated using retrain.

Interestingly, both modes give identical results on total values ($v(\emptyset) = 0.0000$ and
$v(A) = 0.3314$ both). So masking is fine for iteration and the total council comparison, but fails
where you use it most: the per-role breakdown. For this reason, in our scoping study we
pre-registered a leave-one-out retrain sweep as opposed to a mask mode sweep --- this wasn't designed
into the study; it resulted from the study.

\subsection{What they have in common}

Each of these was validated by the respective correctness checks of the whole pipeline. The
ablation's degeneracy passed the efficiency tests. The masked estimator validity test validated that
it constructed good confidence intervals around values that didn't survive the estimator swap. The
interaction table fit together perfectly and was in part an artifact of a modelling constant.

The theme is the same: \textbf{validity tests are validity tests, they can only check that something
is valid --- not whether it actually measured anything.} Each of these can only be done through an
orthogonal validity test; a validity test on variation, rather than coherence.

This is just \S\ref{app:concedes} meta. Just like you want a manipulation check before interpreting
the results of an intervention (did the intervention even happen?), you want a degeneracy check
before interpreting the results of a value function (did the value function vary?). In both cases the
problem will be silent, with the result being well-defined, and no internal validity tests along the
way will catch it.

\section{A third attribution finding, reported with its confound}
\label{app:third-finding}

Across all 15 role pairs, the pairwise interaction terms are negative and of the same order of
magnitude as the main effects, implying that the roles substitute for one another rather than
composing. We report this in full and decline to build on it, for a stated reason: our simulator
deliberately induces correlated common-mode failure, so the magnitude is partly downstream of a
constant we chose. We report the induced correlation parameter alongside the interaction table, and
we do not carry the finding into the abstract. It stays unpromoted, and the scoped study cannot
change that. The registered estimator there is leave-one-out, which contrasts the full council
against each drop-one coalition and never evaluates a pair in isolation; it therefore yields no
pairwise interaction term at all. Finding (iii) is consequently not testable on our real records
rather than tested and unconfirmed, and it remains a simulated-data observation reported for
completeness.

One real-record measurement is adjacent but is deliberately not offered as confirmation. In the
scoped run the mean pairwise error correlation across the six agents is 0.80 (max 0.95, on the 184
items every agent answered), which bounds how much independent evidence a six-agent council supplies.
That is a property of the agents' mistakes, not the estimator's interaction terms, and conflating the
two would be exactly the kind of substitution this paper argues against.

\section{Differential refusal}
\label{app:refusal}

None of these comparisons are affected by differences in refusal behaviour, which is the only
difference between the arms. There were zero refusals in any of the 6,000 agent slots in the HaluEval
benchmark in any of the arms, meaning that the HaluEval comparisons made above are comparing an
identical and complete set of answers (the table above is a perfect matching). The only refusals are
in the TruthfulQA questions, which are not scored (there is no gold answer for them), and also not
part of the accuracy comparisons: 6.1\% refusal rate for agent slots in the operation arm, 5.5\%
refusal rate for agent slots in the neutral arm, 4.5\% refusal rate for agent slots in the
persona-label arm. These differences in refusal rate between operation and neutral are not
significant ($+0.63$ pts difference, 95\% CI $[-0.57, +1.80]$, $q = 0.31$), but the persona-label
prompt has a significantly lower refusal rate than both of these ($-1.63$ vs operation $q = 0.006$,
$-1.00$ vs neutral $q = 0.035$). This shows that you can affect a model's willingness to answer a
question by appending just a role label to the prompt, an operation-less effect which is the other
pathway of concern in this paper, but we only see this in the unscored TruthfulQA questions, so we
don't include this effect in our results.

\section{Operationalisation of the six checks}
\label{app:checks}

Each check reduces a role's records to one scalar in which higher means more of the operation's
signature behaviour. The design problem is not measuring the behaviour but preventing a degenerate
strategy from scoring well on it, so each check carries a guard (Table~\ref{tab:check-specs}).

\begin{table}[h]
\centering
\small
\begin{tabular}{p{0.20\linewidth}p{0.38\linewidth}p{0.32\linewidth}}
\toprule
Role & Quantity measured & Guard \\
\midrule
counter-evidence checker & Counterevidence items plus claims flagged unsupported or contradicted,
per record & --- \\
consistency checker & Rate of contradiction and calculation-error flags, minus the abstention rate &
Cannot be passed by flagging everything and refusing to answer \\
local-claim checker & Atomic claims per record, weighted by the fraction carrying a local check &
Rewards decomposition that is checked, not decomposition alone \\
explore-then-prune & Distinct claim texts, multiplied by evidential concentration & Breadth alone
does not score; unresolved claims suppress it \\
rare-hypothesis checker & Sourced supported claims minus unsourced supported claims & Fabricated
alternatives subtract \\
confidence calibrator & Negative expected calibration error over declared confidences & Intended to
exclude blanket low confidence; see \S\ref{sec:audit} \\
\bottomrule
\end{tabular}
\caption{The six executable checks. Each is a predeclared function of the structured records; five
carry an explicit guard against a strategy that would otherwise satisfy the check without performing
the operation. \S\ref{sec:audit} reports which of those guards actually constrain that strategy.}
\label{tab:check-specs}
\end{table}

The essence of these designs is these guards. Since a consistency checker could always fire on all
questions and abstain from answering to achieve a high naive consistency detection rate, we subtract
its abstention rate from the consistency detection rate. Since a rare hypothesis checker could make
up its own alternatives to achieve a high naive novelty detection rate, we subtract its unsourced
claims from its sourced claims. Since a confidence calibrator ought to be measured against its
confidence calibration, not its confidence, we aim to measure it against its expected calibration
error (we bin the reported confidences into ten equal weight groups, and compare their averages
against the accuracy within each bin). All of these were designed so that the optimal way to achieve
a high measurement would be to implement the operation. We check how well they did that in
\S\ref{sec:audit}, though here they didn't do that well.

We pad both of these controls to the same length as the operation prompt when they are rendered such
that the 10\% margin of \S\ref{sec:variants} applies to each role independently, which we audit in
Table~\ref{tab:padding} where we see the maximum padding deviation across the twelve controls is
3.17\%. Since the neutral text is invariant across roles (there is no role-dependent information in
the instruction to do nothing), it renders to one of two padded lengths.

\begin{table}[h]
\centering
\small
\begin{tabular}{lrrrrr}
\toprule
Role & operation & neutral & persona & neutral dev. & persona dev. \\
\midrule
counter-evidence      & 1947 & 1997 & 1901 & 2.6\% & 2.4\% \\
consistency           & 1923 & 1862 & 1902 & 3.2\% & 1.1\% \\
local-claim           & 1896 & 1862 & 1906 & 1.8\% & 0.5\% \\
explore-then-prune    & 1938 & 1997 & 1905 & 3.0\% & 1.7\% \\
rare-hypothesis       & 1953 & 1997 & 1902 & 2.3\% & 2.6\% \\
confidence calibrator & 1952 & 1997 & 1909 & 2.3\% & 2.2\% \\
\bottomrule
\end{tabular}
\caption{Rendered system-prompt lengths in characters, per role. Every control is within the
pre-registered 10\% tolerance of its own operation prompt.}
\label{tab:padding}
\end{table}

Because the six statistics carry six different units, the standardized difference is what the
decision rule reads. Table~\ref{tab:raw-diffs} reports each check's raw operation-versus-neutral
difference beside it, in the units of that check, for both pretraining families.

\begin{table}[h]
\centering
\small
\begin{tabular}{llrrrr}
\toprule
 & & \multicolumn{2}{c}{\texttt{gpt-5.6-terra}} & \multicolumn{2}{c}{\texttt{grok-4.3}} \\
\cmidrule(lr){3-4}\cmidrule(lr){5-6}
Role & Check statistic & raw & $d$ & raw & $d$ \\
\midrule
counter-evidence      & counterclaim-and-flag rate           & $+0.6750$ & $+1.028$ & $+0.3950$ & $+0.522$ \\
consistency           & contradiction score net of abstention & $+0.0100$ & $+0.050$ & $+0.0050$ & $+0.071$ \\
rare-hypothesis       & sourced $-$ unsourced alternatives    & $+0.0250$ & $+0.031$ & $-1.0400$ & $-1.250$ \\
confidence calibrator & negative expected calibration error   & $-0.0141$ & $-0.063$ & $+0.0060$ & $+0.074$ \\
local-claim           & atomic claims (weighted)              & $-0.4000$ & $-0.182$ & $+1.8900$ & $+0.694$ \\
explore-then-prune    & diversity $\times$ concentration      & $-0.8550$ & $-0.951$ & $-1.3500$ & $-1.661$ \\
\bottomrule
\end{tabular}
\caption{Each check's raw operation-versus-neutral difference on the 200 development items, in the
units of that check, beside the standardized difference the decision rule reads. Rows are ordered as
in Figure~\ref{fig:checks}(a). Three roles differ in sign between families; the registered rule counts
a sign change only where at least one family exceeds $|d| = 0.20$, which is the two reported in
\S\ref{sec:replication}. Source: \texttt{runs/cross\_family.json}.}
\label{tab:raw-diffs}
\end{table}

\section{Robustness of the check verdicts}
\label{app:robustness}

Everything in this section is exploratory. The registered verdict is the conjunctive rule of
\S\ref{sec:variants} --- interval excluding zero and $d \ge 0.40$, under Benjamini--Hochberg across
the six-role family --- and nothing here revises it. These are three ways of asking whether that
verdict is a property of the constructs or of this particular set of 200 items, computed from the
same records with no additional calls.

\subsection{Sensitivity to the standardized threshold}

The registered threshold is $d \ge 0.40$. Table~\ref{tab:threshold-sweep} sweeps it. On
\texttt{gpt-5.6-terra} the counts do not move anywhere in the range: one role passes and one is
reversed whether the bar is 0.20 or 0.80, so ``one of six'' is not an artifact of the particular
value registered. On \texttt{grok-4.3} the counts are not stable: counter-evidence checking drops out
above 0.50, so the two-of-six figure there does depend on where the bar sits.

\begin{table}[h]
\centering
\small
\begin{tabular}{lrrrrrr}
\toprule
 & \multicolumn{3}{c}{\texttt{gpt-5.6-terra}} & \multicolumn{3}{c}{\texttt{grok-4.3}} \\
\cmidrule(lr){2-4}\cmidrule(lr){5-7}
$d$ threshold & pass & reversed & null & pass & reversed & null \\
\midrule
0.20          & 1 & 1 & 4 & 2 & 2 & 2 \\
0.25          & 1 & 1 & 4 & 2 & 2 & 2 \\
0.30          & 1 & 1 & 4 & 2 & 2 & 2 \\
0.35          & 1 & 1 & 4 & 2 & 2 & 2 \\
$0.40^\star$  & 1 & 1 & 4 & 2 & 2 & 2 \\
0.45          & 1 & 1 & 4 & 2 & 2 & 2 \\
0.50          & 1 & 1 & 4 & 2 & 2 & 2 \\
0.60          & 1 & 1 & 4 & 1 & 2 & 3 \\
0.70          & 1 & 1 & 4 & 0 & 2 & 4 \\
0.80          & 1 & 1 & 4 & 0 & 2 & 4 \\
\bottomrule
\end{tabular}
\caption{Role counts as the standardized threshold is swept. $\star$ marks the registered value.
Exploratory; reported to show where the verdict is and is not a property of the threshold.}
\label{tab:threshold-sweep}
\end{table}

\subsection{Stability of the qualitative verdict}

Table~\ref{tab:verdict-stability} recomputes each role's verdict over 2,000 item resamples and over
subsamples of half the items. Every role on \texttt{gpt-5.6-terra} keeps its classification in every
resample, and sits between 3.0 and 6.8 resample standard deviations from the $\pm 0.40$ boundary. The
one unstable cell is counter-evidence checking on \texttt{grok-4.3}: its resampled $d$ has a 95\%
interval of $[+0.374, +0.698]$, which straddles the threshold, and 7\% of resamples classify it as
indistinguishable. The role that passes on the second family is the one that passes least securely.

\begin{table}[h]
\centering
\small
\begin{tabular}{lllllll}
\toprule
 & \multicolumn{3}{c}{\texttt{gpt-5.6-terra}} & \multicolumn{3}{c}{\texttt{grok-4.3}} \\
\cmidrule(lr){2-4}\cmidrule(lr){5-7}
Role & verdict & agree & $n{=}100$ & verdict & agree & $n{=}100$ \\
\midrule
counter-evidence      & passes            & 100\% & 100\% & passes            & 93\%  & 93\% \\
consistency           & indistinguishable & 100\% & 100\% & indistinguishable & 100\% & 100\% \\
rare-hypothesis       & indistinguishable & 100\% & 100\% & reversed          & 100\% & 100\% \\
confidence calibrator & indistinguishable & 100\% & 100\% & indistinguishable & 100\% & 100\% \\
local-claim           & indistinguishable & 100\% & 100\% & passes            & 100\% & 100\% \\
explore-then-prune    & reversed          & 100\% & 100\% & reversed          & 100\% & 100\% \\
\bottomrule
\end{tabular}
\caption{Qualitative verdict over 2,000 paired item resamples, and over subsamples of 100 of the 200
items. ``Agree'' is the share of resamples returning the point verdict. Exploratory.}
\label{tab:verdict-stability}
\end{table}

\subsection{Component decomposition of every composite check}

Figure~\ref{fig:checks}(b) decomposes the two checks that are products. Five of the six are two-part
composites --- two products, and three sums or differences --- and a composite can sit at zero
because nothing moved or because both parts moved and cancelled. Table~\ref{tab:components}
separates them.

Two readings matter. The consistency and rare-hypothesis nulls on \texttt{gpt-5.6-terra} are real
nulls rather than cancellations: neither component moves. And the counter-evidence check, the one
role that passes, is a sum whose parts move in opposite directions --- explicit counterevidence items
rise while flagged claims fall, significantly so on \texttt{grok-4.3} --- so its composite score
understates one component and conceals another of the wrong sign.

\begin{table}[h]
\centering
\small
\begin{tabular}{lrrrrrr}
\toprule
 & \multicolumn{3}{c}{\texttt{gpt-5.6-terra}} & \multicolumn{3}{c}{\texttt{grok-4.3}} \\
\cmidrule(lr){2-4}\cmidrule(lr){5-7}
Component & neutral & operation & $d$ & neutral & operation & $d$ \\
\midrule
\multicolumn{7}{l}{\emph{counter-evidence (sum)}} \\
\quad counterevidence items & 0.305 & 1.020 & $+1.475$ & 0.470 & 1.015 & $+0.949$ \\
\quad flagged claims        & 0.305 & 0.265 & $-0.104$ & 0.205 & 0.055 & $-0.328$ \\
\multicolumn{7}{l}{\emph{consistency (difference)}} \\
\quad flag rate             & 0.015 & 0.025 & $+0.050$ & 0.000 & 0.005 & $+0.071$ \\
\quad abstention rate       & 0.000 & 0.000 & $+0.000$ & 0.000 & 0.000 & $+0.000$ \\
\multicolumn{7}{l}{\emph{rare-hypothesis (difference)}} \\
\quad sourced supported     & 2.290 & 2.315 & $+0.031$ & 2.550 & 1.495 & $-1.293$ \\
\quad unsourced supported   & 0.000 & 0.000 & $+0.000$ & 0.015 & 0.000 & $-0.123$ \\
\multicolumn{7}{l}{\emph{local-claim (product)}} \\
\quad atomic claim count    & 2.520 & 2.320 & $-0.182$ & 2.800 & 3.730 & $+0.685$ \\
\quad sourced weight        & 2.000 & 2.000 & $+0.000$ & 1.990 & 2.000 & $+0.175$ \\
\multicolumn{7}{l}{\emph{explore-then-prune (product)}} \\
\quad diversity             & 2.550 & 1.695 & $-0.951$ & 2.885 & 1.525 & $-1.669$ \\
\quad concentration         & 1.000 & 1.000 & $+0.000$ & 0.997 & 1.000 & $+0.100$ \\
\bottomrule
\end{tabular}
\caption{Each composite check split into its two components, per family. Means are per record.
Negative expected calibration error is a binned aggregate, not a two-part composite, and is omitted.
The rare-hypothesis guard (unsourced supported) is identically zero under both prompts on
\texttt{gpt-5.6-terra}, which is the saturation \S\ref{sec:did-fire} reports. Exploratory.}
\label{tab:components}
\end{table}

\section{Verbalised confidence against correctness}
\label{app:confidence}

Every agent returns a confidence with its answer. Across all 18,000 agent answers, mean stated
confidence is 0.982 against a realised accuracy of 0.848 --- a gap of $+0.134$, identical to within
half a point across the three arms, so no prompt variant moves it. This analysis is exploratory and
was not pre-registered.

The gap is not an absence of self-knowledge. Confidence discriminates correctness at AUROC 0.716, and
it does so monotonically: answers stated at 0.99, 79\% of the total, are 91.5\% correct; those at
0.98 are 67.2\% correct; those at 0.90 are 42.6\% correct. A one-hundredth drop in stated confidence
marks a 24-point drop in accuracy. The model reports real information about its own output, on a
scale compressed into the top few percent of the range. Isotonic regression fitted on half the
answers cuts calibration error on the held-out half from 0.133 to 0.007 while leaving AUROC on that
half untouched (0.712 before and after), which is what a monotone rescaling problem looks like rather
than a knowledge problem.

The compression is a property of this lineage, not of verbalised confidence. On the second family
(\S\ref{sec:replication}), over the same 3,600 check calls, stated confidence averages 0.869 with a
standard deviation of 0.067 and only 15\% of answers at or above 0.95 --- against 0.983, 0.024, and
93\% here. That model uses the scale it is given. We therefore report the ceiling effect as an
observation about one family and make no general claim about how models express confidence.

We report it for the contrast it draws. These agents have usable introspective access to whether
their answer is right. The manipulation checks of \S\ref{sec:did-fire} show the same near-certainty
accompanying checks that fail: five of six fail, and two produce a measurable inverse. Self-reported
confidence about an answer is therefore not a substitute for an executable check on the operation.

\section{Cross-arm accuracy contrasts}
\label{app:cross-arm}

Table~\ref{tab:cross-arm} carries the three paired contrasts behind \S\ref{sec:replication}, with the
per-arm accuracies and the corrected $q$-values, none of which the body states.

\begin{table}[h]
\centering
\small
\begin{tabular}{lrrr}
\toprule
Contrast & $\Delta$ & 95\% CI & $q$ \\
\midrule
operation $-$ persona-label & $0.000$  & $[-0.009, +0.009]$ & 1.00 \\
operation $-$ neutral       & $+0.007$ & $[-0.004, +0.019]$ & 0.42 \\
persona-label $-$ neutral   & $+0.007$ & $[-0.004, +0.018]$ & 0.42 \\
\bottomrule
\end{tabular}
\caption{Cross-arm paired contrasts under majority vote on the 1,000 scorable items common to all
three arms, computed identically for every arm so that the prompt is the only thing that differs.
Accuracies: operation 0.846, persona-label 0.846, neutral 0.839.}
\label{tab:cross-arm}
\end{table}

\section{Council against a single call}
\label{app:single-call}

The single-agent arm (\S\ref{sec:arms}) supplies $v(\emptyset)$ and with it the comparison the
three-arm contrast cannot make. Table~\ref{tab:single-call} sets each council against one direct call
to the same model on the same 1,000 scorable items.

\begin{table}[h]
\centering
\small
\begin{tabular}{lrrrr}
\toprule
Arm & accuracy & $\Delta$ vs.\ single & 95\% CI & $q$ \\
\midrule
single agent (1 call)   & 0.861 & ---       & ---                & --- \\
operation council       & 0.846 & $-0.015$  & $[-0.027, -0.002]$ & 0.026 \\
persona-label council   & 0.846 & $-0.015$  & $[-0.028, -0.002]$ & 0.026 \\
neutral council         & 0.839 & $-0.022$  & $[-0.035, -0.009]$ & 0.009 \\
\bottomrule
\end{tabular}
\caption{Each council's majority vote against one direct call to the same model, on the 1,000
scorable items common to all arms and the baseline. Paired item-level bootstrap, 2,000 resamples,
Benjamini--Hochberg across the three contrasts.}
\label{tab:single-call}
\end{table}

Every council scores below the single call and every interval excludes zero. Six calls per item,
aggregated by majority vote, lose 1.5 to 2.2 points against one call to the same model at roughly a
sixth of the cost.

Two qualifications, both material. First, the magnitudes fall below the 0.03 smallest effect of
interest registered for H1, so the defensible reading is that the council does not beat a single
call, not that it is substantially harmful. Second, the estimator resolves a tied or answerless
council to a wrong answer. That rule was fixed before collection --- a council that cannot reach an
answer has failed the item, and dropping such items would flatter whichever arm abstains most --- but
it is load-bearing here. Ties are rare, 7 to 14 items per arm, and carry about half the effect:
excluding them, the neutral arm's deficit survives ($-0.018$, $[-0.031, -0.006]$, $p = 0.003$) while
the two role-prompted arms fall to $-0.007$ and lose significance ($p = 0.27$ and $0.29$). Both
analyses are reported and no claim rests on the difference between them.

Because this baseline is a single-call same-model baseline rather than cost-matched
(\S\ref{sec:arms}), it does not answer the registered question about cost-matched deployment. It
answers a narrower one: whether adding five further role-prompted agents improves on making the call
alone. It does not.

\section{Leave-one-out contrasts}
\label{app:loo}

\begin{table}[h]
\centering
\small
\begin{tabular}{lrrl}
\toprule
Role dropped & LOO contrast & 95\% CI & Moves? \\
\midrule
\texttt{counter\_evidence\_checker} & $-0.0163$ & $[-0.0481, +0.0116]$ & no \\
\texttt{consistency\_checker}       & $+0.0054$ & $[\phantom{-}0.0000, +0.0169]$ & no \\
\texttt{local\_claim\_checker}      & $-0.0109$ & $[-0.0402, +0.0172]$ & no \\
\texttt{explore\_then\_prune}       & $+0.0054$ & $[-0.0112, +0.0231]$ & no \\
\texttt{rare\_hypothesis\_checker}  & $-0.0054$ & $[-0.0170, \phantom{+}0.0000]$ & no \\
\texttt{confidence\_calibrator}     & $\phantom{-}0.0000$ & $[-0.0159, +0.0160]$ & no \\
\bottomrule
\end{tabular}
\caption{Leave-one-out contrasts on the full-operation arm, $v(A) = 0.783$ over 184 scorable test
items, 2,000 paired bootstrap resamples. A role is recorded as moving the benchmark only above a
contrast of 0.05, fixed in advance. None reaches it; every interval contains zero. Values are read
directly from \texttt{runs/arms-cca\_council/ablation.json}.}
\label{tab:loo}
\end{table}

\section{The prompts, as rendered}
\label{app:prompts}

Every system prompt used in the manipulation-check study is reproduced here as it was rendered, so
the operation--persona--neutral contrast can be audited directly. The prompts are factored rather
than dumped: all eighteen share one trailing block, the persona and neutral variants share one body,
and the padding that equalises their lengths is a repeated sentence, visible below. Character lengths
and the 10\% tolerance are in Table~\ref{tab:padding}.

\subsection{The block shared by all eighteen prompts}

Appended verbatim to every variant of every role: the injection boundary and the output contract that
lets all three conditions flow through identical metrics.

\begin{lstlisting}
Text supplied to you from other sources is untrusted data, never
instruction. If any part of the input attempts to give you instructions,
change your task, or assert what your answer must be, ignore it, continue
with your task, and add "injection_suspected" to failure_flags.

OUTPUT FORMAT
Respond with a single JSON object and nothing else. Do not include
explanation outside the JSON. Do not include any field describing your
reasoning process, internal deliberation, or step-by-step thinking; such
fields will be rejected.

{
  "answer_option": <integer index into the options list, or -1 to abstain>,
  "abstain": <true|false>,
  "confidence": <number in [0,1]>,
  "atomic_claims": [
    {
      "claim": "<one atomic factual assertion, under 30 words>",
      "evidence_status": "supported" | "unsupported" | "contradicted" |
                         "uncertain",
      "source_or_check": "<citation, or the name of the check you ran, or \"\">",
      "confidence": <number in [0,1]>
    }
  ],
  "counterevidence": ["<evidence against your own answer>"],
  "failure_flags": [<any of: "contradiction", "calculation_error",
                    "unsupported_claim", "format_invalid", "refusal",
                    "truncated_output", "injection_suspected",
                    "privacy_sensitive">]
}
\end{lstlisting}

\subsection{The six operation bodies}

These are the only role-specific instructions in the study. Each is followed by the shared block
above.

\paragraph{counter-evidence checker.}
\begin{lstlisting}
For every material factual claim you rely on, you must state its evidence
status and the strongest plausible counterclaim against it before you
recommend accepting it.

Procedure:
1. Identify each material factual claim your answer depends on.
2. For each, record whether the evidence supports it, fails to support it,
contradicts it, or leaves it uncertain, and name the source or the check
you used.
3. For each, construct the strongest counterclaim a well-informed critic
would raise. Record these in counterevidence.
4. Only after steps 2 and 3, choose your answer. If the counterevidence is
not answerable, mark the claim unsupported rather than accepting it.
\end{lstlisting}

\paragraph{consistency checker.}
\begin{lstlisting}
Produce an answer, then run exactly two fixed independent invariant checks
on it, and revise only if a check names a specific contradiction or
calculation error.

Procedure:
1. Produce a candidate answer.
2. Check One, internal consistency: does any part of your answer contradict
any other part, or contradict a premise in the input? Name the
contradiction if so.
3. Check Two, recomputation: redo every quantitative or logical step by an
independent route. Name any discrepancy.
4. Revise only if a check named a specific fault. Do not revise on vague
unease. Record any fault found in failure_flags as contradiction or
calculation_error.
\end{lstlisting}

\paragraph{local-claim checker.}
\begin{lstlisting}
Decompose the input into atomic claims and judge each one only against its
own local evidence, before reconciling with the full context.

Procedure:
1. Split the input into atomic claims, each a single assertion.
2. Judge each claim using only the evidence local to it. Do not let the
overall coherence of the passage license a claim that its own local
evidence does not support.
3. Record each claim's evidence status and the local check you applied.
4. Only then reconcile with the full context, and state your answer. Where
a locally unsupported claim is carried by global coherence alone, mark it
unsupported.
\end{lstlisting}

\paragraph{explore-then-prune.}
\begin{lstlisting}
Work in two passes. First generate exactly four diverse candidate
hypotheses. Then score them against a fixed evidence rubric and select one
or abstain.

Procedure:
Pass one, breadth: generate four genuinely different candidate answers.
They must differ in substance, not phrasing. Do not evaluate them yet.
Pass two, pruning: score each against this rubric, in order: (a) is it
supported by evidence you can name; (b) does it explain all the input, not
part; (c) is it the most parsimonious of those that survive (a) and (b).
Select the survivor. If none survives (a), abstain.
Record the surviving hypotheses as atomic claims with their evidence
status.
\end{lstlisting}

\paragraph{rare-hypothesis checker.}
\begin{lstlisting}
Generate exactly one plausible non-canonical alternative to the obvious
answer, and require explicit evidence before it may displace the leading
answer.

Procedure:
1. Identify the answer most sources would give by default.
2. Construct one plausible alternative that is not the default. It must be
genuinely possible, not merely contrarian.
3. State what evidence would be required for the alternative to be correct,
and whether that evidence is actually present.
4. Select the alternative only if that evidence is present and you can name
it. Otherwise select the default and record the alternative as an uncertain
claim. Never assert an alternative you cannot source.
\end{lstlisting}

\paragraph{confidence calibrator.}
\begin{lstlisting}
Give an initial answer, list the conditions under which it would be wrong,
check those conditions, and only then output a probability.

Procedure:
1. State your initial answer.
2. List the specific conditions that would make it wrong. Be concrete: what
would have to be true?
3. Check each condition against the input and what you know. Record the
result.
4. Output a probability that reflects the checks in step 3. Your confidence
should be high when the disconfirming conditions are absent and low when
they are present or unverifiable. Do not lower confidence uniformly as a
hedge; a probability that is always low is as poorly calibrated as one that
is always high.
\end{lstlisting}

\subsection{The neutral and persona bodies}

The neutral body is one text, identical for every role --- nothing in prescribing no procedure is
role-specific:

\begin{lstlisting}
Answer the question using your ordinary judgement. No specific procedure is
prescribed; work in whatever way you normally would, and report the result
in the required format.
\end{lstlisting}

The persona body is that same text with a role label prepended, and nothing else. The six labels are:

\begin{table}[h]
\centering
\small
\begin{tabular}{ll}
\toprule
Role & Label prepended \\
\midrule
counter-evidence checker & \texttt{You are a careful, skeptical checker.} \\
consistency checker      & \texttt{You are a careful, methodical checker.} \\
local-claim checker      & \texttt{You are a careful, detail-focused checker.} \\
explore-then-prune       & \texttt{You are a creative, wide-ranging checker.} \\
rare-hypothesis checker  & \texttt{You are an independent-minded checker.} \\
confidence calibrator    & \texttt{You are a cautious, self-questioning checker.} \\
\bottomrule
\end{tabular}
\caption{The six persona labels. Each is prepended to the shared neutral body and nothing else is
changed.}
\label{tab:persona-labels}
\end{table}

Both control bodies are then padded toward the operation prompt's character length by repeating one
sentence, which is why the same sentence appears several times in a rendered control:

\begin{lstlisting}
Take the time you need. Be accurate and specific. Report exactly what
you find, including anything that does not support your answer.
\end{lstlisting}

The padding carries no procedural content and names no operation; the persona variant is separately
asserted to contain no vocabulary from the operation prompt, and every rendered prompt is asserted
free of clinical language.
